\section{Discussion: what the discrepancies have in common}
\label{sec:discussion}

Four of the six claims did not survive measurement, and it would be possible
to stop there and call the result a list of corrections. That is the least
useful reading available. In three of the four the underlying argument is
sound; what fails is the step from the argument to the quantity reported
beside it. This section makes that step explicit, because it is the part that
generalises past the two architectures at hand.

\subsection{The counting asymmetry, worked out}

Section~\ref{sec:complexity} deferred the reason a fifty-fold ratio measures
as a two-fold one. It can be recovered from the reported figure itself.

Write $M$ for the token count, $d$ for the channel width, and $N$ for the
state width of a selective scan. The self-attention figure that appears in
these comparisons is the whole block: the quadratic term $2M^2d$ for the
attention matrices, plus projection terms of order $4Md^2$ for the queries,
keys, values, and output. The state space figure set beside it is $8MDN$ per
scan direction, with $D$ the expanded inner width. That is the scan kernel
and nothing else.

But a selective scan block is not a scan. It has an input projection, an
output projection, and two further projections that produce its
input-dependent parameters, and all of these scale as $Md^2$ exactly as
attention's do. At the operating point of the table in question they are an
order of magnitude larger than the scan term they surround, because $N$ is a
few tens while $d$ is a few hundreds. One side of the comparison counted its
projections and the other did not.

Counting both sides whole, on the same models and the same device, gives the
ratio in Table~\ref{tab:e1}. The cross-over survives the correction; only its
size does not.

\subsection{A claim can be true on one axis and measured on another}

Put three of the discrepancies side by side and they have the same shape.

The complexity argument---linear in tokens against quadratic---is true of the
scan kernel, and the number quoted from it was read as though it were true of
the block. The anisotropy claim is true of the field near its peak, where the
scan geometry acts, and the covariance index that was used to test it is
dominated by the broad skirt. The dilution argument is true of
normalisation, and the quantity it was taken to predict, how much attribution
lands on the object, depends also on the width of the recurrent state and on
whether the object's tokens are contiguous in the scan.

In none of the three would ``the claim is false'' be right, and in none of
them would ``the claim is true'' be honest. The accurate statement names the
axis: the claim holds on the axis its argument is about, and the measurement
was taken on a different one. Our verdicts are phrased that way for this
reason, and not to be diplomatic.

The fourth discrepancy is a different kind and we keep it separate. The
throughput and memory figures are a scope mismatch rather than an axis
mismatch: same quantity, different resolution, different device, different
memory budget. The two call for different repairs. An axis mismatch is
repaired by restating the claim; a scope mismatch is repaired by widening the
measurement, which we have not done and do not claim to have done.

\subsection{Accuracy and cost do not tell the same story}

The factorial experiment leaves the accuracy interaction undetermined in size
and clear in sign. Had we summarised the experiment on the accuracy axis
alone---as architecture comparisons ordinarily are summarised---the honest
report would have been that structure and operator barely interact. The
largest interaction in the experiment would have gone unmentioned, because it
is in wall-clock cost and not in top-1.

This is worth stating as more than an anecdote. The two axes are not
redundant, and which one a comparison is summarised on is a choice that
changes its conclusion.

Two limits on what the experiment separates should be read alongside it. Our
\emph{structure} factor bundles hierarchical downsampling with convolutional
locality, and nothing here says which of the two carries the effect;
separating them is a $2 \times 2 \times 2$ design and three times the
compute. And at sixty-four tokens neither operator is under the pressure that
distinguishes them---the sequence is short enough that the choice of mixer has
little room to matter, which is itself one of the findings and not a defect of
the design.

\subsection{What pre-registration and tooling did here}

Two registered predictions failed, in Section~\ref{sec:dilution} and
Section~\ref{sec:factorial}, and both failures produced something. The first
sent us to check whether the finding depended on the scale it was measured
on; it did, which is why we report the result at the level rather than as a
slope. The second sent us to inspect the validation split rather than accept
a pleasant number. Neither check would have been prompted if the predictions
had been written after the measurements, and both changed what we report.

The fused-kernel undercount described in Section~\ref{sec:complexity} is the
tooling counterpart of the same discipline. It matters beyond this paper
because it is not symmetric: the most convenient way to produce a FLOP
comparison between these two families understates both models and understates
the state space model considerably more, and nothing in the output says so.
A number that is wrong in a known direction is more dangerous than one that
is merely noisy.

None of this shows that the arguments in the literature are wrong. It shows
that the arguments and the numbers printed next to them are, in several
well-known cases, about different things.
